\documentclass[11pt,a4paper]{article}
\usepackage[utf8]{inputenc}
\usepackage[T1]{fontenc}
\usepackage[spanish]{babel}
\usepackage{hyperref}
\usepackage{geometry}
\geometry{margin=2.5cm}
\usepackage{parskip}

\title{075 el truncado se lee como completo}
\author{Alejandro Valenzuela Rivera\\\small Universidad Autónoma de Nuevo León}
\date{\today}
\begin{document}
\maketitle

\begin{quote}
\textbf{Sesgo declarado.} Este trabajo se sostiene desde el sesgo declarado de su autor, y no finge neutralidad. Lo que afirma se puede tumbar, y ese es el punto: el falsificador está escrito, no escondido.
\end{quote}

\section*{Resumen}
Un mensaje que no cabe en el canal no es claro: es truncado. Y el truncado se lee como completo. Eso pasa con quien no entiende el idioma, quien no oye, quien no ve y quien procesa distinto.

\section{Por qué se sostiene}
El que emite cree que comunico, el que recibe cree que entendio, y el error no tiene donde verse. Y como parece entero, nadie reclama. Por eso Lengua y Acomodo son las dos letras de ALAS que cargan con la accesibilidad, y Sostenido impide que la ayuda se vuelva dependencia: una persona que depende del que le explica no fue comunicada, fue asistida.

\section{De dónde viene}
ALAS y la conversacion del 2026-10-25.

\section{El caso}
Ninguna: es el caso real de la accesibilidad.

\section{Método}
Un mensaje que dice que se trunco, cuando se trunco.

\section{Límites, y qué lo tumba}
Un mensaje truncado que se presenta como completo.

\section{Lo que queda abierto}
\textquestiondown{}Que no cabe en el canal de esta persona, y quien lo declara?

\section*{Cómo se leyó esta tesis}
Esto se escribió en el repositorio man-made-empathy, del que este candado es una pieza: candados/06 del caso/075 el truncado se lee como completo.md. La máquina que lo compuso lo hizo desde este mismo archivo, y no agregó nada que no estuviera escrito. Si algo de esto es falso, se dice en el mismo sitio, y se corrige con fecha.

\end{document}
